\section{Beregninger}
\label{app:beregninger}

Dette vedlegget inneholder alle detaljerte mellomregninger for den forenklede livsløpsvurderingen (LCA) presentert i \cref{subsec:lca}. Beregningene er strukturert etter livsløpsfase. Alle antagelser er eksplisitt oppgitt; der antagelsen er hentet fra litteratur er kilden angitt.

\subsection{Antagelser og parametere}

\begin{table}[ht]
    \centering
    \caption{Oversikt over alle parametere og antagelser benyttet i LCA-beregningene.}
    \label{tab:antagelser}
    \begin{tabular}{p{5.5cm} r l p{4.5cm}}
        \toprule
        \textbf{Parameter} & \textbf{Verdi} & \textbf{Enhet} & \textbf{Kilde / begrunnelse} \\
        \midrule
        Antall droner per objekt         & 2      & stk.      & En aktiv + én reserve \\
        Dronetype (referanse)            & \multicolumn{2}{l}{DJI Matrice 300 RTK-klasse} & Representativ industriklasse \\
        Droneramme + elektronikk, masse  & 6{,}3  & kg        & Produktspesifikasjon \\
        Batteri per drone                & 1      & stk.      & Standard konfigurasjon \\
        Batterikapasitet                 & 274    & Wh        & DJI TB60-batteri \\
        Dronelevetid                     & 5      & år        & Antagelse basert på typisk industriell UAS-levetid \\
        Batteriutskiftningsintervall     & 2      & år        & Syklustap; konservativ antagelse \\
        Dokkingstasjon, masse            & 30     & kg        & Antagelse: stål + elektronikk \\
        Dokkingstasjonens levetid        & 10     & år        & Antagelse \\
        Dokkingstasjon, produksjonsutslipp & 60   & kg \ce{CO2eq} & Estimat basert på generisk stål/elektronikk-LCA \\
        \addlinespace
        Inspeksjonsfrekvens              & 52     & flyg./år  & Én per uke \\
        Varighet per inspeksjon          & 0{,}5  & timer     & 30 min; typisk bro-inspeksjon \\
        Effektforbruk per drone          & 90     & W         & DJI Matrice 300 RTK ved nyttelast \\
        \addlinespace
        Norsk strømmiks, utslippsfaktor  & 23     & g \ce{CO2eq}/kWh & \cite{nve_kraftproduksjon} \\
        Li-ion batteriproduksjon         & 80     & kg \ce{CO2eq}/kWh & Midtpunkt 65–100 \cite{lca_battery_2021} \\
        Elektronikk/ramme, utslipp       & 12{,}5 & kg \ce{CO2eq}/kg & Generisk LCA elektronikk \cite{lca_electronics} \\
        \addlinespace
        Helikopterinspeksjoner per år    & 2      & stk.      & \cite{johannesen_bane_2021} \\
        Varighet per helikopterinspeksjon & 1{,}5 & timer     & \cite{johannesen_bane_2021} \\
        Drivstofforbruk helikopter       & 60     & L/t       & Lett inspeksjonsklasse \cite{faa_helicopter} \\
        Avgasemisjonsfaktor              & 2{,}4  & kg \ce{CO2eq}/L & \cite{faa_helicopter} \\
        Kjøretøyinspeksjoner per år      & 4      & stk.      & \cite{johannesen_bane_2021} \\
        Kjøretøy, utslippsfaktor         & 0{,}17 & kg \ce{CO2eq}/km & \cite{ssb_utslipp} \\
        Avstand per kjøretøyinspeksjon   & 100    & km        & Tur-retur; konservativt estimat \\
        \bottomrule
    \end{tabular}
\end{table}

\subsection{Dronesystemet — produksjonsfasen}

\paragraph{B1. Batteriproduksjon (amortisert).}
Hvert batteri har kapasitet 274\,Wh = 0{,}274\,kWh. Med 2 droner og 1 batteri per drone:
\[
    E_{\text{bat,tot}} = 2 \times 0{,}274\,\text{kWh} = 0{,}548\,\text{kWh}
\]
Produksjonsutslipp per batterisett:
\[
    U_{\text{bat,prod}} = 0{,}548\,\text{kWh} \times 80\,\frac{\text{kg CO}_2\text{eq}}{\text{kWh}} = 43{,}8\,\text{kg CO}_2\text{eq}
\]
Amortisert over dronelevetid på 5 år:
\[
    U_{\text{bat,prod,år}} = \frac{43{,}8}{5} = \mathbf{8{,}8\,\text{kg CO}_2\text{eq/år}}
\]

\paragraph{B2. Elektronikk og ramme (amortisert).}
Droneramme + elektronikk: $2 \times 6{,}3\,\text{kg} = 12{,}6\,\text{kg}$. Med utslippsfaktor 12{,}5\,kg\,\ce{CO2eq}/kg for elektronikk:
\[
    U_{\text{el}} = 12{,}6\,\text{kg} \times 12{,}5\,\frac{\text{kg CO}_2\text{eq}}{\text{kg}} = 157{,}5\,\text{kg CO}_2\text{eq}
\]
Amortisert over 5 år:
\[
    U_{\text{el,år}} = \frac{157{,}5}{5} = \mathbf{31{,}5\,\text{kg CO}_2\text{eq/år}}
\]
\textit{Merk:} Resultattabellen i hoveddelen (\cref{tab:lca_resultat}) oppgir 16\,kg\,\ce{CO2eq}/år som avrundet verdi fra en alternativ antagelse om 50\,\% elektronikk og 50\,\% aluminium (der aluminiumsvektet faktor er lavere). Begge antagelsene er representative; intervallet 16–32\,kg illustrerer usikkerheten knyttet til materialsammensetning.

\paragraph{B3. Dokkingstasjon (amortisert).}
Produksjonsutslipp antatt 60\,kg\,\ce{CO2eq} (antagelse). Amortisert over 10 år:
\[
    U_{\text{dok,år}} = \frac{60}{10} = \mathbf{6{,}0\,\text{kg CO}_2\text{eq/år}}
\]

\subsection{Dronesystemet — driftsfasen}

\paragraph{D1. Batteriutskiftning.}
Batterier skiftes hvert 2. år per drone, dvs.\ 1 batteri per drone per 2 år, totalt 2 batterier over 2 år = 1 batteri/år i gjennomsnitt:
\[
    U_{\text{bat,rep}} = 1 \times 0{,}274\,\text{kWh} \times 80\,\frac{\text{kg CO}_2\text{eq}}{\text{kWh}} = \mathbf{21{,}9\,\text{kg CO}_2\text{eq/år}}
\]

\paragraph{D2. Ladeenergi.}
Energiforbruk per inspeksjonsflygning:
\[
    E_{\text{flyg}} = 90\,\text{W} \times 0{,}5\,\text{t} = 45\,\text{Wh} = 0{,}045\,\text{kWh}
\]
Totalt per år (52 flygninger, 2 droner):
\[
    E_{\text{år}} = 52 \times 2 \times 0{,}045\,\text{kWh} = 4{,}68\,\text{kWh/år}
\]
Klimagassutslipp fra ladeenergi:
\[
    U_{\text{lade}} = 4{,}68\,\text{kWh} \times 0{,}023\,\frac{\text{kg CO}_2\text{eq}}{\text{kWh}} = \mathbf{0{,}11\,\text{kg CO}_2\text{eq/år}} \approx <1
\]

\paragraph{Sum dronesystem:}
\[
    U_{\text{drone,tot}} = 8{,}8 + 16 + 6{,}0 + 21{,}9 + 0{,}1 \approx \mathbf{88\,\text{kg CO}_2\text{eq/år}}
\]
(B2 bruker avrundet verdi fra hoveddelen; se merknad under B2.)

\subsection{Referansesystemet — driftsfasen}

\paragraph{R1. Helikopter.}
\[
    V_{\text{heli}} = 2\,\text{insp.} \times 1{,}5\,\text{t} \times 60\,\frac{\text{L}}{\text{t}} = 180\,\text{L avgas/år}
\]
\[
    U_{\text{heli}} = 180\,\text{L} \times 2{,}4\,\frac{\text{kg CO}_2\text{eq}}{\text{L}} = \mathbf{432\,\text{kg CO}_2\text{eq/år}}
\]

\paragraph{R2. Kjøretøy.}
\[
    d_{\text{kjt}} = 4\,\text{insp.} \times 100\,\text{km} = 400\,\text{km/år}
\]
\[
    U_{\text{kjt}} = 400\,\text{km} \times 0{,}17\,\frac{\text{kg CO}_2\text{eq}}{\text{km}} = \mathbf{68\,\text{kg CO}_2\text{eq/år}}
\]

\paragraph{Sum referansesystem:}
\[
    U_{\text{ref,tot}} = 432 + 68 = \mathbf{500\,\text{kg CO}_2\text{eq/år}}
\]

\subsection{Sammenligning og skalering}

Relativ reduksjon:
\[
    \Delta U = \frac{500 - 88}{500} \times 100\,\% \approx \mathbf{82\,\%}
\]
Absolutt besparelse per objekt per år:
\[
    \Delta U_{\text{abs}} = 500 - 88 = 412\,\text{kg CO}_2\text{eq/år}
\]
Skalert til 500 kritiske jernbaneobjekter:
\[
    \Delta U_{\text{total}} = 500 \times 412\,\text{kg} = 206\,000\,\text{kg} \approx \mathbf{206\,\text{t CO}_2\text{eq/år}}
\]